% concatenated from 250107311.tex
%#!pdflatex -synctex=1
\documentclass[a4paper,12pt]{article}
%
\usepackage{amssymb,amsmath}
\usepackage[legacycolonsymbols]{mathtools}
\usepackage{bm}%
\usepackage{ascmac}
\usepackage{mathrsfs}
\usepackage{stmaryrd}
\usepackage{comment}
\usepackage[pdftex]{hyperref}
\hypersetup{
	colorlinks=true,
	citecolor=blue,
	linkcolor=blue,
	urlcolor=orange,
}
\usepackage{arydshln}
%\usepackage{multirow,bigdelim}
\usepackage{bigdelim}
\usepackage{physics}
\usepackage{cite}
\usepackage{comment}

\setlength{\textheight}{23cm}
\setlength{\textwidth}{16cm}
\setlength{\topmargin}{0cm}
\setlength{\headheight}{0pt}
\setlength{\oddsidemargin}{0pt}
\setlength{\evensidemargin}{0pt}
%\setlength{\unitlength}{0.7mm}
%
\def\nn{\nonumber}
\def\Z{\mathbb{Z}}
%\def\Z3{\mathbb{Z}_2^3}
%\def\Z4{\mathbb{Z}_2^4}
\def\Zn{\mathbb{Z}_2^n}
%\def\bH{\bar{H}}
%\def\bE{\bar{E}}
\def\bH{\mathsf{h}}
\def\bE{\mathsf{e}}
\def\DP#1#2{\hat{#1}\cdot \hat{#2}}
\def\PH#1#2{(-1)^{\hat{#1}\cdot \hat{#2}}}
\def\alg{$\Z2$-SuperPoincar\'e algebra}
\def\affine{$\mathbb{Z}_2^2$-$\widehat{sl(2)}$}
\def\cI{\mathcal{I}}
\def\cIh{\hat{\mathcal{I}}}
\def\sF#1{\mathscr{F}_{#1}}
\def\PBa#1#2{\{#1(y), #2(x) \}}
\def\cT{\mathcal{T}}
\def\cU{\mathcal{U}}
\def\ua{\alpha}
\def\da{\dot{\alpha}}
\def\ub{\beta}
\def\db{\dot{\beta}}
\def\g{\mathfrak{g}}
\def\h{\mathfrak{h}}
\def\HLS{Haag-\L opusza\'nski-Sohnius }
\def\CM{Coleman-Mandula }
\def\ev{\text{ev.}}
\def\od{\text{od.}}

%
%\renewcommand{\thesection}{\Roman{section}.}
\renewcommand{\theequation}{\thesection.\arabic{equation}}
%\renewcommand{\theequation}{\arabic{section}.\arabic{equation}}
%\renewcommand{\thefigure}{\thesection.\arabic{figure}}
%
%\newtheorem{lemma}{Lemma}[section]
%\newtheorem{prop}[lemma]{Proposition}
%\newtheorem{thm}[lemma]{Theorem}
%\newtheorem{DEF}[lemma]{Definition}
%\newtheorem{cor}[lemma]{Corollary}
%
%\renewcommand{\thelemma}{\arabic{section}.\arabic{lemma}}

%
\usepackage{amsthm}
\theoremstyle{plain}
\newtheorem{thm}{Theorem}
\newtheorem*{thm*}{Theorem}

\theoremstyle{definition}
\newtheorem{dfn}{Definition}
\usepackage{version}	% required for `\comment' (yatex added)
\begin{document}

\title{Novel possible  symmetries of $S$-matrix generated by $\Z_2^n$-graded Lie superalgebras}


\author{Ren Ito and Akio Nago
	\\[15pt]
	 Department of Physics, Osaka Metropolitan University,
\\Sugimoto Campus, Osaka 558-8585, Japan}


\maketitle
%\begin{flushright}
%	\today
%\end{flushright}
%\thispagestyle{empty}

\vfill
\begin{abstract}
In this paper, we explore the $\Z_2^n$-graded Lie (super)algebras as novel possible generators of symmetries of $S$-matrix. 
As the results, 
we demonstrate that 
a $\Z_2^n$-graded extension of the supersymmetric algebra can be a symmetry of $S$-matrix. Furthermore, it turns out that a $\Z_2^n$-graded Lie algebra appears as internal symmetries.
They are natural extensions of \CM theorem and \HLS theorem, which are the no-go theorems for generators of symmetries of $S$-matrix. 
%It is known a generator of supersymmetry can avoid the no-go theorem, called Coleman-Mandula theorem, on the symmetry allowed for the $S$-matrix in quantum field theory. In addition, it is also known the supersymmtry can be extended by the $\Z_2^n$-grading ($\Z_2^n$-supersymmetry).
%However it has not been known that the $\Z_2^n$-supersymmetry is allowed for the $S$-matrix in quantum field theory yet. 
%In this paper we show that a generator of $\mathbb{Z}_2^n$-supersymmetry can avoid the Coleman-Mandula theorem in the same to the supersymmetry.
\end{abstract}



%%%%%%%%%%%%%%%%%%%%%%%%%%%%%%%%%%%%%%%%%%%%%%%%%%%%%%%%%%%%%%%%%%%%%%%%%%%%%%%%%%%%%%

%\clearpage
%\setcounter{page}{1}

%\tableofcontents
\section{Introduction}
Since the introduction of $SU(4)$ symmetry in the classification of nuclei by Wigner in 1937, the properties and fundamental structure of particles have been described using group theory\cite{F.J.D, S.W}.
The quark model introduced in the Standard Model is a model using $SU(3)$ flavor symmetry proposed by Gell-Mann in 1964\cite{M.G}. 
It successfully describes the relationship among hadrons with the same spin.
Furthermore, if we consider the $SU(6)$ symmetry for quark flavor and spin, we can embed  hadrons with different spins in the same representation. 

However, such $SU(6)$ symmetry that unify and describe different spins are approximate symmetries that exist in the non-relativistic limit, and despite some success, all attempts to generalize them to relativistic quantum theory have failed\cite{S.C}. 
Indeed, it was shown by Coleman and Mandula in 1967 that such symmetries are forbidden in relativistic quantum field theory (QFT), known as the Coleman-Mandula theorem\cite{CM}.
The theorem says that all possible symmetries of $S$-matrix generated by Lie algebra are only Poincar\'e symmetry and internal symmetry such that its generators commute with Poincar\'e algebra.
%This theorem is too powerful to build a theory on.
This prevents us to unify Poincar\'e symmetry and internal symmetries into a single multiplet.
%This is because symmetry is one of the most important guiding principles for theory in physics, and particularly the Lie algebra generated by the Lie group describes continuous symmetry.
%So, it is natural to consider some ways to avoid the Coleman-Mandula theorem.
The reason for the strong restriction of this theorem is that the symmetries allowed in the S-matrix of QFT were searched only in the framework of Lie algebras. 
If we want to avoid the \CM theorem, it is natural to explore symmetries generated by  algebras beyond the Lie algebra.

In searching for symmetries that do not conflict with the Coleman-Mandula theorem, Haag, \L opsuza\'nski and Sohnius extended the algebra to Lie superalgebra, so that the symmetry generators are closed not only  for comutation relations but also for  anti-commutation relations\cite{HLS}.
This showed that supersymmetry is the only possible symmetry as an extension of Poincare symmetry, which is known as the Haag-\L opsuza\'nski-Sohnius theorem.
As its result, if we consider the Lie superalgebra as symmetry generators of the $S$-matrix, supersymmetry is allowed in QFT. 

%Moreover, supersymmetry has some nice properties like the relation to the solvability of system. Therefore, many supersymmetric models were considered in many parts of physics. 
%%%%%%%%%%%%%%%%%%%%%%%%%%%%%%%%%%%%%%%%%%%%%%%%%%%%%%%%%%%%%%%
%%%%%%%%%%%%%%%%%%%%%%%%%%%%%%%%%%%%%%%%%%%%%%%%%%%%%%%%%%%%%%%
In 1978,  further extensions of the Lie algebra were introduced by Rittenberg and Wyler \cite{RW1,RW2}(see also \cite{Ree,sch1}). This is called $\Z_2^n$-graded Lie (super)algebra, which is the extended algebra with the grading by abelian group $\Z_2^n:= \Z_2\times \Z_2 \times \cdots $(repeated $n$ times). It reduces to an ordinary Lie algebra when $n=0$,  and to a Lie superalgebra when $n=1$.
  
This type of algebra exhibits more complicated (anti-)commutation relations and possesses interesting properties in both physics and mathematics. In physics, symmetries generated by $\Z_2^n$-graded Lie superalgebras have been identified in various contexts, such as de Sitter SUGRA \cite{Vas}, nuclear quasi-spin \cite{jyw},  parastatistics \cite{tol},  non-relativistic  Dirac  equation \cite{AKTT1,AKTT2}, etc. 
Additionally, several $\Z_2^n$-graded extension of spacetime supersymmetries have  also been proposed \cite{zhe,LR,Toro1,Toro2,tol2,Bruce}. In particular, $\Z_2^n$-graded extension of super-Poincar\'e algebra introduced by Bruce \cite{Bruce} has some interesting properties and $\Z_2^n$-graded version of classical and quantum mechanics,  sigma model, etc. have been discussed based on the extension \cite{BruDup,AAD,AAd2,AKTcl,AKTqu,DoiAi1,DoiAi2,AiDoi,brusigma,bruSG,AIKT,Topp,Topp2,AiItoTa,AiItoTa2}. 

This indicates that $\Z_2^n$-graded Lie (super)algebras are not uncommon in physics. Instead, they provide new insights into the analysis of physical systems based on symmetries, which is one of the key motivations for studying their structure and representations. 
The structure and representation theories of $\Z_2^n$-graded Lie (super)algebras have been continuously explored since their introduction \cite{sch3,SchZha,Sil,ChSiVO,SigSil,PionSil,CART,MohSal,NAJS,NAPIJS,StoVDJ,Meyer,IsStvdJ,NAPSIJSinf,NAKA,KTclassification,Meyer2,NeliJoris,NeliJoris2,LuTan,FaFaJ,AiSe}.

However, despite the extensive studies carried out so far, the application of $\Z_2^n$-graded Lie (super)algebras to $(1+3)$-dimensional relativistic quantum field theory has not been considered at all. To this end, one of the most important issues is the $\Z_2^n$-graded Lie (super)algebra can be symmetries of $S$-matrix. 
If this is the case, it provides an extension of the \HLS theorem.
%A major reason for this is their applicability to relativistic quantum field theory.
%This is because, while the applicability to field theory has been investigated for Lie algebras and Lie superalgebras through the \CM theorem and the , no such investigation has been carried out for \( \mathbb{Z}_2^n \)-graded Lie (super)algebras regarding these restrictions. 
The purpose of this paper is to investigate the possible \( \mathbb{Z}_2^n \)-graded Lie (super)algebra symmetries of $S$-matrix in relativistic quantum field theory, following the approach of the \CM theorem and \HLS theorem. 
As the results, 
we demonstrate that 
a $\Z_2^n$-graded extension of the supersymmetric algebra can be a symmetry of $S$-matrix. Furthermore, it turns out that a $\Z_2^n$-graded Lie algebra appears as internal symmetries.
%They are natural extensions of \CM theorem and \HLS theorem, which are the no-go theorems for generators of symmetries of $S$-matrix. 

Our paper is organized as follows. In \S 2, we recall the definition of $\Z_2^2$-graded Lie superalgebra\cite{RW1,RW2}, and  briefly review \CM theorem and \HLS theorem\cite{CM,HLS} in the case of massive particles. 
In \S 3, we consider the novel generators of symmetries of $S$-matrix in the case of massive particles. We first analyze the scenario with $\Z_2^2$-grading, following the frameworks of \CM theorem and \HLS theorem. Subsequently, using the same method as for the $\Z_2^2$-grading, we will consider the possible generators with $\Z_2^n$-gradings. Finally we will discuss the relation between the generators of internal symmetries and $\Z_2^n$-graded Lie algebra.



%%%%%%%%%%%%%%%%%%%%%%%%%%%%%%%%%%%%%%%%%%%%%%%%%%%%%%%%%%%%%%%%%%%%%%%%%%%%%%%%%%%%%

\section{Preliminaries}
\subsection{$\Z_2^n$-graded Lie superalgebra}
%The purpose of this paper is to investigate novel possible generators of symmetries of $S$-matrix within the framework of $\Z_2^n$-graded Lie superalgebras(LSAs). 
We recall the definition of the $\Z_2^n$-graded LSA\cite{RW1,RW2}.
$\Z_2^n$ refers to the $n$-fold direct product of the abelian group $\mathbb{Z}_2=\{0,1\}$, namely $\Z_2^n=\mathbb{Z}_2\times \mathbb{Z}_2 \times \cdots$ (repeated $n$ times). 
\begin{dfn}
 Let $\g$ be a $\Z_2^n$-graded vector space  (over $\mathbb{R}$ or $ \mathbb{C}$), which is the direct sum of homogeneous vector subspaces labeled by an element of $\Z_2^n$:
\begin{align}
 \g = \bigoplus_{\vec{a}\in\Z_2^n}\g_{\vec{a}}.\label{Z2n-VS}
\end{align}
An element of $ \g_{\vec{a}}$ is said to have the $\Z_2^n$-grading $\vec{a}$. 
\end{dfn}
\begin{dfn}\label{defsalg}
We define the $\Z_2^n$-Lie bracket by
\begin{equation}
	\llbracket A, B \rrbracket = AB - (-1)^{\vec{a}\cdot\vec{b}} BA, \label{Z2n-LB}
\end{equation}
where $A \in \g_{\vec{a}}, \ B \in \g_{\vec{b}}$, and $ \vec{a}\cdot\vec{b} $ is the standard scalar product of $n$-dimensional vectors. 
A $\Z_2^n$-graded vector space is called a $\Z_2^2$-graded Lie superalgebra if 
$ \llbracket A,B \rrbracket \in \g_{\vec{a}+\vec{b}} $ and the Jacobi identity is satisfied:
\begin{align}
 \llbracket A, \llbracket B,C \rrbracket \rrbracket
= \llbracket \llbracket A,B \rrbracket, C \rrbracket
+ (-1)^{\vec{a}\cdot \vec{b}} \llbracket B, \llbracket A, C \rrbracket \rrbracket. \label{Z2n-JI}
\end{align}
\end{dfn}

The $\Z_2^n$-graded LSA can be regarded as an extended Lie algebra because the $\Z_2^n$-Lie bracket is a commutator (anti-commutator) for $ \vec{a}\cdot\vec{b} $ is even (odd). Actually when $n=0$, a $\Z_2^n$-graded LSA is equivalent to an ordinary Lie algebras. Furthermore, when $n=1$, it corresponds to the standard LSA, which consists of a bosonic sector ($0$-grading) and a fermionic sector ($1$-grading). It is known that for $n\ge 2$, a $\Z_2^n$-graded LSA also generates a symmetry. 
%%%%%%%%%%%%%%%%%%%%%%%%%%%%%%%%%%%%%%%%%%%%
\subsection{Coleman-Mandula theorem}
%%%%%%%%%%%%%%%%%%%%%%%%%%%%%%%%%%%%%%%%%%%%
In this subsection, we briefly review the Coleman-Mandula theorem, which is known as the no-go theorem for symmetries in QFT \cite{J.M, CM}.
The theorem states that
``The continious symmetries of the $S$-matrix in QFT can only be generated by a Lie algebra consisting of the translation generator $P_{\mu}$, the homogeneous Lorentz transformation generator $M_{\mu \nu}$, and the internal symmetry generator $B_l$''. 
Namely they discussed possible symmetries of the $S$-matrix generated by $\Z_2^n$-graded LSA's with $n=0$ and proved Poincar\'e symmetry and internal symmetry were only allowed in QFT. 
To consider the avoidance of this theorem, we look at the proof of this theorem\cite{J.M, S.W}.
They used the following resonable assumptions in their proof,
\begin{enumerate}
\item
The $S$-matrix is defined in relativistic local field theory, 
\item
The finite number of particle types with masses lighter than an arbitrary mass $M > 0$ is finite, 
\item
The amplitudes for elastic two-body scattering are analytic functions of the scattering angle at almost all energies and angles. 
\end{enumerate}
The symmetry generator here means a Hermitian operator on the Hilbert space which is commutative with the $S$-matrix and acts on multi particle states as well as single particle ones.
Note that we work in the momentum space so that any matrix element of the group generator is a distribution in momentum space.

%They proved the no-go theorem on the symmetry allowed for the $S$-matrix in quantum field theory. They used the following assumptions in their proof of this theorem,
%\begin{enumerate}
%\item
%The S matrix is based on relativistic local field theory. 
%\item
%The variety of particles with masses lighter than an arbitrary mass $M > 0$ is finite. 
%\item
%There is an energy gap between the vacuum and the single particle state. 
%\end{enumerate}
%According to the theorem, the only possible Lie algebra consists of the generators translations by $P_{\mu}$, the generator of homogeneous Lorentz transformations by $M_{\mu \nu}$, and the generators of internal symmetries by $B^l$.  
%The set of 
First, we consider all generators which commute with $P_{\mu}$.
%we call $\mathfrak{C}$. 
Any compact Lie algebra can be decomposed into a direct sum of a semi-simple part $\mathfrak{g}^\mathcal{S}$ and an Abelian part $\mathfrak{g}^\mathcal{A}$ (Levi-Decomposition). 
%When $\mathfrak{C}$ is decomposed into two parts, 
Since $P_\mu$ commute with itself, $P_{\mu}$ is in Abelian part: $P_{\mu} \in \mathfrak{g}^\mathcal{A}$. 
Thus, we consider the generators $B_{\alpha} \in \mathfrak{g}^\mathcal{S} $. 
We transform the commutation relation between $P_{\mu}$ and $B_\alpha$ by Lorentz transformation
\begin{align}
[U(\Lambda) B_\alpha U^{-1}(\Lambda),{\Lambda^\nu}_\mu P_\nu]=0
\end{align}
where $U(\Lambda)$ is the unitary operator on Hilbert space and represents the Lorentz transformation.
Since $ {\Lambda^\nu}_\mu $ is regular, $U(\Lambda) B_\alpha U^{-1}(\Lambda)$ must be a linear combination of $B_\alpha$:
\begin{align}
U(\Lambda) B_\alpha U^{-1}(\Lambda) = {D(\Lambda)^\beta}_\alpha B_\beta,
\end{align}
where $D(\Lambda)$ is the representation matrix of Lorentz group.
Since the metric of semi-simple part $\mathfrak{g}^\mathcal{S}$ is positive definite, $D(\Lambda)$ must be an identity matrix for $B_\alpha$ to satisfy the same algebraic relation after Lorentz transformation.
Noting that the Lorentz group is not compact, the symmetry generated by $B_\alpha \in \mathfrak{g}^\mathcal{S} $ is independent of Lorentz symmetry.
The $B_\alpha$ is independent of the momentum and acts as an identity matrix with respect to the spin, which is the generator of internal symmetry $B_l$.
%In this theorem, they considered all generators of $G$ that are commutative with $P_{\mu} $. 
%Let us call the set of these $\mathfrak{c}$. 
%The generators $B^{\alpha} \in \mathfrak{c} $ that compose the Lie algebra are commutative with $M_{\mu \nu}$.
%Then the generator $B^{\alpha}$ are independent of spin and momentum, which is only an internal symmetry $B^l$.

Next, we consider all symmetry generators which do not commute with $P_\mu$.
All symmetry generators $G_{\ua}$ act on single particle states as follows
\begin{align} \label{alloperator}
G_\alpha \ket{n ,p}= \sum_{n^\prime} \int \dd^4 p^{\prime} \left( K_\alpha(p^{\prime}, p) \right)_{n^{\prime}, n}\ket{n^\prime ,p^\prime} ,
\end{align}
where $n^{\prime}, n$ denote the spin and particle type.
The kernel $K_\alpha (p^{\prime}, p)$ will be zero if both $p$ and $p^{\prime} $ are not on the mass hyperboloid. 
Furthermore, the kernel $K(p^{\prime}, p)$ is zero for all $p \neq p^\prime$:
\begin{align}
K_\alpha(p^{\prime}, p) = K_\alpha(p) \delta^4(p-p^{\prime}) .
\end{align}
From the second assumption, since the mass of the particle takes only discrete values, $G_\alpha$ commute with the mass square operator.
Thus, the commutation relation of $P_\mu$ and $G_\alpha$ satisfies
\begin{align}
[P^\mu,G_\alpha]=a_\alpha^{\mu\nu}P_\nu,
\end{align}
where $a_\alpha^{\mu \nu}$ is anti-symmetric with respect to $\mu,\nu$.
It follows that $G_\alpha$, which does not commute with $P_\mu$, is the Lorentz generator 
%This relationship of the part of $G_\alpha$ which do not comute with $P_\mu$ shows exactly the Poincar\' e symmetry:
\begin{align}\label{Poincare}
[P_\mu, M_{\nu\rho}]=i(g_{\nu\mu}P_\rho - g_{\rho \mu}P_\nu).
\end{align}
Therefore, the symmetries of the $S$-matrix in QFT are only the Poincar\'e and internal symmetries.
%\begin{align}
%G_\alpha = \frac{i}{2}a^{\nu\mu}_\alpha M_{\mu\nu} + B_\alpha
%\end{align}

%%%%%%%%%%%%%%%%%%%%%%%%%%%%%%%%%%%%%%%%%%%%%%%%%%%%%%%%%%%%%%%%%%%%%%%%%%%%%%%%%%%%%%%%%%%%%%%%%%%%%%%%%%%%

\subsection{Haag-\L opusza\'nski-Sohnius theorem}
In this subsection, we briefly review the Haag-\L opusza\'nski-Sohnius theorem\cite{HLS}, which extends the Coleman-Mandula theorem.
In \eqref{alloperator}, $G_\alpha$ can be separated into two different types; bosonic type for same statistics $n^{\prime}, n$, or fermionic type for different statistics $n^{\prime}, n$.
All bosonic type generators are fully analyzed in the Coleman-Mandula theorem.
A complete basis of bosonic type generators are given by the $P_{\mu}$,$M_{\mu\nu}$ and  $B^l$.
On the other hand, the fermionic type generator is not mentioned in the Coleman-Mandula theorem because of the assumption that the all generators consist of Lie algebra.
They extended the Lie algebra to include anti-commutation relations in order to consider the fermionic type generators.
%This algebra is called a LSA, specifically a $\Z_2^n$-graded LSA with $n=1$.
This theorem says that,
``Poincar\'e symmetry can only be extended to supersymmetry''

Now, let us see how this supersymmetry avoids the Coleman-Mandula theorem.
Note that $G_\alpha$ in \eqref{Poincare} can be bosonic or fermionic. 
However, the proof in the previous section concludes that the possible $G_\alpha$ is only \eqref{Poincare}.
Therefore, it is enough to consider fermionic generator $Q$ commute with $P_\mu$. 
%Since the Coleman-Mandula theorem restricts the part which do not commutes with $P_\mu$ to generators of general symmetry, we consider all fermionic generators $Q$, which commute with $P_\mu$.
Let $\mathfrak{c}$ denote the set of all generators commute with $P_\mu$.
%In order to see specific $Q$, we need only look at  the representation of the Lorenz group as in the proof of the previous theorem.
Since $\mathfrak{c}$ is stable under the Lorentz transformations, it is a finite dimensional representation space of Lorentz groups, and all generators in $\mathfrak{c}$ can be decomposed into irreducible representation.
The irreducible representation of the Lorentz group is specified by two spinors, which we write as $(a,b)$-rep .
In general, because fermionic generators can have an odd number of spin indices, we denote the generator of  $(j,j^\prime)$-rep  and its conjugate $(j^\prime,j)$-rep  characterizing two independent spins $\ua,\db$ as $Q_{\ua_1\dots\ua_{2j}\db_1\dots\db_{2j'}}, \ \bar{Q}_{\ub_1\dots\ub_{2j^\prime}\da_1\dots\da_{2j}}$,
%Let $\mathfrak{c}$ be the set whose elements commute with the $P_\mu$. 
%From the conclusin (B), there are only a finite number of linearly indepemdemt $[11],[01],[10]$ type generators in $\mathfrak{c}$ since $K(\rho)$ has finite dimensionality and G is fixed by specifying two such matrices. Moreover $\mathfrak{c}$ is stable under homogeneous Lorentz transformations.
%Therefore this set is a representation space of LG.
%Irrep of elements of LG is labeled by pair of indices $(j,j)$ each integer or half integer, i.e. 
%$Q_{\ua_1\ua_2,...,\ua_{2j}\da_1\da_2,...,\da_{2j'}}$, which is symmetric in the $2j$ undotted and the $2j'$ dotted indices. Let us call irrep $(j,j')$. Due to the spin-stat. thm, $2(j+j')$ must be odd for fermi type generators.
where $j,\ j^\prime$ is an integer or half-integer, and $2j$ undotted and $2j^\prime$ dotted indices are symmetric in each part.
Now consider the anti-commutation relation between $Q$ and $\bar{Q}$.
Since it is producted from two fermionic generators, it must be bosonic and belong to $(j+j^\prime, j+j^\prime)$-rep.
The bosonic generators are defined by the Coleman-Mandula theorem and are either $P_\mu,M_{\mu\nu},B_l$. The irreducible representation of the Lorenz group of each generator is as follows
\begin{align}
P_\mu \rightarrow \left(\frac12,\frac12 \right)\text{-rep},\quad M_{\mu\nu} \rightarrow \left(1,0\right)\text{-rep} +\left(0,1\right)\text{-rep},\quad B_l \rightarrow \left(0,0\right)\text{-rep}. \nn
\end{align}
It follows that $j+j^\prime \leq \frac{1}{2}$, which means $Q \in \mathfrak{c}$ is only rank-1 spinor charge. 
Therefore, the anti-commutator of $Q$ with $\bar{Q}$ becomes $P_{\mu}$
\footnote{
Our notations are as follows. Let $g_{\mu\nu}=\mathrm{diag}(+,-,-,-)$ and $(\sigma^\mu)_{\ua\da}=(\sigma^0,\sigma^1,\sigma^2,\sigma^3)_{\ua\da}$, where
\begin{align*}
 \sigma^0=\smqty(\pmat{0}),\qquad
\sigma^1= \smqty(\pmat{1})
,\qquad
\sigma^2= \smqty(\pmat{2})
,\qquad
\sigma^3= \smqty(\pmat{3}).
\end{align*}
%The simbol ``bar'' means ``dagger'' in quantum theory. Then $(\bar \sigma^\mu)_{\da\ua} =(\sigma^\mu)_{\ua\da}$ and 
Define the complex conjugate for any matrices as $(M_{\ua\ub})^\ast=:\bar M_{\da\db}$ and the Hermitian conjugate as $(M_{\ua\ub})^\dag =:\bar M_{\db\da}$. Then we obtain $\{(\sigma^\mu)_{\ua\db}\}^\dag= (\bar \sigma^\mu)_{\ub\da}=(\sigma^\mu)_{\ub\da}$ and $(\bar \sigma^\mu)^{\da\ua}:=\varepsilon^{\da\db}\varepsilon^{\ua\ub}\sigma^\mu_{\ub\db}=(\sigma^0,-\sigma^i)$, where $\varepsilon_{\ua\ub}$ is the antisymmetric tensor characterized by $\varepsilon^{12}=\varepsilon_{21}=1$(the same for dotted indices). 
}
\begin{align}
\{ Q_{\ua}, \bar{Q}_{\db} \}= \sigma^\mu_{\ua\db}P_{\mu}.
\end{align}
%%%
$Q_\alpha, \bar Q_{\da}$ are the supersymmetry generators.
The other relations are obtained in the same way
\begin{align}
[Q_\ua, M_{\mu \nu}]= {(\sigma_{\mu\nu})_\alpha}^{\ua_1} Q_{\ua_1}, \quad
%-i(\varepsilon_{\ua\ua_1}Q_{\ua_2}+\varepsilon_{\ua\ua_2}Q_{\ua_1}) \\
\{Q_\ua,Q_\ub\}=0.
\end{align}

%The theorem say that the fermionic generator in $\mathfrak{c}$ is only rank-1 spinorial charge $Q_{\ua}, \bar{Q}_{\da}$ which occur supersymmetry.  
If we consider that there are several charges $Q_{\ua\, L}, \bar{Q}_{\da\, M}\ (L,M=1,2,\dots )$, anti-commutation relations of $Q_{\ua, L}$ is non-zero:
\begin{align}
\{ Q_{\ua\, K}, Q_{\ub\, L}\} 
&= \epsilon_{\ua \ub} \sum_l a_{l,KL}B_{l} \eqqcolon \epsilon_{\ua \ub} Z_{KL},\nn\\
\{\bar Q_{\da\, K}, \bar Q_{\db\, L}\} 
&= \varepsilon_{\da \db} \sum_l \bar a_{l,KL}B_{l} \eqqcolon \varepsilon_{\da \db} \bar Z_{KL},
%= \epsilon_{\ua \ub} Z_{LM} .
\end{align}
where $a^l_{LM}$ is antisymmetric with respect to $L,M$.
From the Jocobi identity, we can get
\begin{align}\label{central}
[ Z_{KL}, Z_{MN}] = [ Z_{KL}, \bar Z_{MN}] = 0.
%, \quad [Z_{KL},  Q_{\ua, K}]=0 .
\end{align}
Since $Z_{KL}$ commutes with itself, $Z_{KL}$ is in the Abelian part $\mathfrak{g}^{\mathcal{A}}$ of the Lie algebra and commutative with $B_l \in \mathfrak{g}^\mathcal{S}$.
Therefore, $Z_{KL}$ commutes with all elements of the algebra. 
It is called the central charge.
The LSA $\langle P_{\mu}, M_{\mu\nu}, Z_{KL}, \bar Z_{KL}, Q_{\ua\, L}, \bar{Q}_{\da\, M}\rangle$ is called the superPoincar\'e algebra, and the LSA $\langle P_{\mu}, Z_{KL}, \bar Z_{KL}, Q_{\ua\, L}, \bar{Q}_{\da\, M}\rangle$ is called the supersymmetric algebra.
%Besides the fact that the Lie algebra $\mathcal{A}$ of $B^l$ is compact, it can be decomposed into a direct sum of a semi-simple part $\mathcal{A}_1$ and an Abelian part $\mathcal{A}_2$.
%Thus $ Z_{LM}$ is an invariant subalgebra in $\mathcal{A}$, and also commutative with $G$. 
%$P_{\mu}, B^{l}, Z_{LM}$ and ${\ua, L}$ is Lie superalgebra.
%%%%%%%%%%%%%%%%%%%%%%%%%%%%%%%%%%%%%%%%%%%%%%%%%%%%%%%%%%%%%%%%%%%%%%%%%

%%%%%%%%%%%%%%%%%%%%%%%%%%%%%%%%%%%%%%%%%%%%%%%%%%%%%%%%%%%%%%%%%%%%%%%%%
\section{Novel possible generators of the $S$-matrix}
%%%%%%%%%%%%%%%%%%%%%%%%%%%%%%%%%%%%%%%%%%%%%%%%%%%%%%%%%%%%%%%%%%%%%%%%%
\setcounter{equation}{0}
\subsection{Possible $\Z_2^2$-graded generators of symmetries}
%%%%%%%%%%%%%%%%%%%%%%%%%%%%%%%%%%%%%%%%%%%%%%%%%%%%%%%%%%%%%%%%%%%%%%%%%
In this section, we search for novel possible generators of symmetries of $S$-matrix within the framework of the $\Z_2^2$-graded LSA.
From \eqref{Z2n-VS} $\Z_2^2$-graded LSA $\g$ has four types of homogeneous vector subspaces:
\begin{align}
 \g=\g_{(0,0)}\oplus\g_{(1,1)}\oplus\g_{(0,1)}\oplus\g_{(1,0)},
\end{align}
and, from \eqref{Z2n-LB}, $\Z_2^2$-Lie brackets between each subspace are given, in terms of (anti-)commutator as in Table \ref{Z22-LB}.
\begin{table}[h]
 \begin{align*}
  \begin{array}{c|cccc}
   & \g_{(0,0)}& \g_{(1,1)}& \g_{(0,1)}& \g_{(1,0)}\\ \hline
   \g_{(0,0)}& [,]& [,]& [,]& [,]\\
   \g_{(1,1)}& [,]& [,]& \{,\}& \{,\}\\ 
   \g_{(0,1)}& [,]& \{,\}& \{,\}& [,]\\
   \g_{(1,0)}& [,]& \{,\}& [,]& \{,\}\\
  \end{array}
 \end{align*}
\caption{$\Z_2^2$-Lie brackets between each subspace}
\label{Z22-LB}
\end{table}

\noindent Thus, the generator $G_{\ua}$ in \eqref{alloperator} is categorized into two types of bosonic generators in $\g_{(0,0)}$ and $\g_{(1,1)}$, and two types of fermionic generator in $\g_{(0,1)}$ and $\g_{(1,0)}$. The $(0,0)$-graded generators are (ordinary) bosonic because the algebraic relations between them and all generators are given by the commutator. 
For the $(0,1)$ and the $(1,0)$-graded generators, the relations among generators with the same grading are given by the anti-commutators, while those between generators with different gradings are given by commutators. They are called commuting fermionic generators.
The (1,1)-graded generators are referred to as exotic bosonic ones because the relations among them are given by commutators, whereas their relations with fermionic generators are given by anti-commutators. 
Note that there are four types of symmetry generators of $S$-matrix.
%The types of $G_{\ua}$ are determined by the particle type $n$ and $n'$ in \eqref{alloperator}. 
This means there are four types of particles in QFT within this framework. %The relations between spin and statistics of these particles are shown in the appendix \ref{ap1}.

We focus on the following three subalgebras:
\begin{align*}
 \h_1=\g_{00}\oplus\g_{11},\qquad  \h_2=\g_{00}\oplus\g_{01},\qquad \h_3=\g_{00}\oplus\g_{10}.
\end{align*}
As indicated in Table \ref{Z22-LB}, $\h_1$ is an ordinary Lie algebra, while $\h_2$ and $\h_3$ are nothing but LSA's. 
This shows that we can obtain generators by applying the \CM theorem to $\h_1$ and the \HLS theorem to $\h_2$ and $\h_3$. 
Therefore, in $\h_1$, there are only the translation generators $P_\mu$, Lorentz generator $M_{\mu\nu}$ and two types of generators of internal symmetries: $B_l^{00}\in \g_{(0,0)}$, $B_l^{11}\in\g_{(1,1)}$. 
$P_\mu$ and $M_{\mu\nu}$, of course, commute with $B^{00}_l$ and $B^{11}_l$.
The algebraic relations between $B$'s are given by
\begin{alignat}{2}
 [B_l^{00},B_m^{00}]&=i \sum_k f_{lmk} B_k^{00},\qquad &
 [B_l^{11},B_m^{11}]&=i \sum_k g_{lmk} B_k^{00},\nn\\
 [B_l^{00},B_m^{11}]&=i \sum_k h_{lmk} B_k^{11},
% [B_l^{11},B_m^{00}]&=-ih_{lmk} B_k^{11},
\end{alignat}
where $f_{lmk}$, $g_{lmk}$ must be antisymmetric $f_{lmk}=-f_{mlk}$, $g_{lmk}=-g_{mlk}$, but $h_{lmk}$ has no conditions because $B^{00}_l$ has the different grading from $B^{11}_l$.

On the other hand, $\h_2$ is identical to the LSA introduced in \S 2.3.
%In $\h_2$, we can only add 
Denoting the $(0,1)$-graded supercharges by $Q_{\ua\,K}^{01}$, $\bar Q_{\da\,K}^{01}\in \g_{(0,1)}$,
the algebraic relations between $P_\mu$, $B^{00}_l$ and $Q^{01}$'s are given by
\begin{alignat}{2}
 \{Q^{01}_{\ua\,K}, \bar Q^{01}_{\da\,L}\} &=\sigma^{\mu}_{\ua\da}P_\mu\delta_{KL},\qquad&
 [Q^{01}_{\ua\,K},B_l^{00}]&=\sum_L s_{l,KL}Q^{01}_{\ua\,L}, \nn\\
\{Q^{01}_{\ua\,K}, Q^{01}_{\ub\,L}\} &= \varepsilon_{\ua\ub} \sum_l a_{l,KL} B^{00}_l,\qquad&
\{\bar Q^{01}_{\da\,K}, \bar Q^{01}_{\db\,L}\} &= \varepsilon_{\da\db} \sum_l \bar a_{l,KL} B^{00}_l, 
\end{alignat}
where $a_{l,KL}=-a_{l,LK}$. Introducing the generator $Z^{(1)}_{KL}:=\sum_l a_{l,KL} B^{00}_l\in\g_{(0,0)}$, it becomes the central charge in $\h_2$, which means
\begin{alignat}{2}
\{Q^{01}_{\ua\,K}, Q^{01}_{\ub\,L}\} &= \varepsilon_{\ua\ub} Z^{(1)}_{KL},\qquad&
 [Z^{(1)}_{KL},\,\circ\,]&=0,
\nn\\
\{\bar Q^{01}_{\da\,K}, \bar Q^{01}_{\db\,L}\} &= \varepsilon_{\da\db} \bar Z^{(1)}_{KL},\qquad&
 [\bar Z^{(1)}_{KL},\,\circ\,]&=0.
\end{alignat}
Similarly, $\h_3$ is also identical to the LSA introduced in \S 2.3.
Denoting the $(1,0)$-graded supercharges by $Q_{\ua\,K}^{10}$, $\bar Q_{\da\,K}^{10}\in \g_{(1,0)}$ in $\h_3$, the algebraic relations between $P_\mu$ and $Q^{10}$'s are given by
\begin{alignat}{2}
 \{Q^{10}_{\ua\,K}, \bar Q^{10}_{\da\,L}\} &=\sigma^{\mu}_{\ua\da}P_\mu\delta_{KL},\qquad&
 [Q^{10}_{\ua\,K},B_l^{00}]&=\sum_L t_{l,KL}Q^{10}_{\ua\,L},\nn\\
\{Q^{10}_{\ua\,K}, Q^{10}_{\ub\,L}\} &= \varepsilon_{\ua\ub} \sum_l b_{l,KL} B^{00}_l,\qquad&
\{\bar Q^{10}_{\da\,K}, \bar Q^{10}_{\db\,L}\} &= \varepsilon_{\da\db} \sum_l \bar b_{l,KL} B^{00}_l, 
\end{alignat}
where $b_{l,KL}=-b_{l,LK}$. Introducing the generator $Z^{(2)}_{KL}:=\sum_l b_{l,KL} B^{00}_l \in\g_{(0,0)}$, it also becomes the central charge in $\h_3$, which means
\begin{alignat}{2}
\{Q^{10}_{\ua\,K}, Q^{10}_{\ub\,L}\} &= \varepsilon_{\ua\ub} Z^{(2)}_{KL},\qquad&
 [Z^{(2)}_{KL},\,\circ\,]&=0,
\nn\\
\{\bar Q^{10}_{\da\,K}, \bar Q^{10}_{\db\,L}\} &= \varepsilon_{\da\db} \bar Z^{(2)}_{KL},\qquad&
 [\bar Z^{(2)}_{KL},\,\circ\,]&=0.
\end{alignat}

%$Z^{00}_{KL}$ is the central charge in $\h_2$, and $\widetilde Z^{00}_{KL}$ is the central charge in $\h_3$. If both of $Z^{00}_{KL}$ and $\widetilde Z^{00}_{KL}$ are also the central charges in $\g$, they should be $Z^{00}_{KL}=\widetilde Z^{00}_{KL}$. This implies  $a_{KL}^l=b_{KL}^l$, which shows the number of the $(0,1)$-graded supercharges equals the number of the $(1,0)$-graded supercharges.

We need to determine the other algebraic relations. Taking into account that $Q^{01}_{\ua\,L}$ and $Q^{10}_{\ua\,L}$ are a $(\frac{1}{2},0)$-rep, $[Q^{01}_{\ua\,K},Q^{10}_{\ub\,L}]\in \g_{(1,1)}$ will be a $(1,0)$-rep or $(0,0)$-rep, and  $[Q^{01}_{\ua\,K},\bar Q^{10}_{\da\,L}]\in \g_{(1,1)}$ will be a $(\frac{1}{2},\frac{1}{2})$-rep. Moreover, because $B^{11}_l \in \g_{11}$ is $(0,0)$-rep, we can determine the following relations: 
\begin{alignat}{2}
 [Q^{01}_{\ua\,K},Q^{10}_{\ub\,L}]&=\sum_l \varepsilon_{\ua\ub} c_{l,KL} B_l^{11},\qquad&
 [\bar Q^{01}_{\da\,K},\bar Q^{10}_{\db\,L}]&=- \sum_l \varepsilon_{\da\db} \bar c_{l,KL} B_l^{11}
,\nn\\
 \{Q^{01}_{\ua\,K},B_l^{11}\}&=\sum_L u_{l,KL}Q^{10}_{\ua\,L},\qquad &
 \{Q^{10}_{\ua\,K},B_l^{11}\}&=\sum_L v_{l,KL}Q^{01}_{\ua\,L}
,\nn\\
 [Q^{01}_{\ua\,K},\bar Q^{10}_{\da\,L}]&=0.
\end{alignat}

We compute $\Z_2^2$-Jacobi identities \eqref{Z2n-JI} to consider the conditions for structure constants such that $\g$ becomes $\Z_2^2$-graded LSA. 
We obtain the following conditions
%%%%%%%%%%%%%%%%%%%%%%%%%%%%%%%%%%%%%%%%%%%%%%%%%%%%%%%%%%%%BQQ^dag and QQQ^\dag part
from $(A,B,C)=(Q,\bar Q, B)$:
\begin{alignat}{3}
 s_{l,LK}=&\bar s_{l,KL},\qquad&  t_{l,LK}=&\bar t_{l,KL},\qquad&  u_{l,LK}=&\bar v_{l,KL},
\end{alignat}
from $(A,B,C)=(Q,Q,\bar Q)$:
\begin{alignat}{2}
\sum_l a_{l,KL}\bar s_{l,KL}=&0
,\qquad& 
\sum_l b_{l,KL}\bar t_{l,KL}=&0
,\nn\\ 
\sum_l c_{l,KL}\bar u_{l,KL}=&0
,\qquad& 
\sum_l c_{l,KL}\bar v_{l,KL}=&0
,\label{JfQQQd}
\end{alignat}
%%%%%%%%%%%%%%%%%%%%%%%%%%%%%%%%%%%%%%%%%%%%%%%%%%%%%%%%%%%%QQQ part
from $(A,B,C)=(Q,Q,Q)$:
\begin{alignat}{2}
\sum_{l} \big( a_{l,KL}s_{l,NM} + a_{l,NL}s_{l,KM} \big) &= 0
,\qquad&
\sum_{l} \big( b_{l,KL}t_{l,NM} + b_{l,NL}t_{l,KM} \big) &= 0
,\nn\\
\sum_{l} \big( c_{l,KL}u_{l,NM} + c_{l,NL}u_{l,KM} \big) &= 0
,\qquad&
\sum_{l} \big( c_{l,LK}v_{l,NM} + c_{l,LN}v_{l,KM} \big) &= 0
,\nn\\
\sum_{l} \big( a_{l,KL}t_{l,NM} + c_{l,LN}u_{l,KM} \big) &= 0
,\qquad&
\sum_{l} \big( b_{l,KL}s_{l,NM} - c_{l,NL}v_{l,KM} \big) &= 0
,\label{JfQQQ}
\end{alignat}
and from $(A,B,C)=(Q,Q,B)$:
\begin{align}
 \sum_M \big( s_{l,LM}a_{m,KM}- s_{l,KM}a_{m,LM} \big) &= i\sum_k f_{klm}a_{k,KL}
,\nn\\
 \sum_M \big( u_{l,LM}c_{m,KM}- u_{l,KM}c_{m,LM} \big) &= i\sum_k h_{klm}a_{k,KL}
,\nn\\
 \sum_M \big( t_{l,LM}b_{m,KM}- t_{l,KM}b_{m,LM} \big) &= i\sum_k f_{klm}b_{k,KL}
,\nn\\
 \sum_M \big( v_{l,LM}c_{m,MK}- v_{l,KM}c_{m,ML} \big) &= -i\sum_k h_{klm}b_{k,KL}
,\nn\\
 \sum_M \big( t_{l,LM}c_{m,KM}- s_{l,KM}c_{m,ML} \big) &= -i\sum_k g_{lkm}c_{k,KL}
,\nn\\
 \sum_M \big( v_{l,LM}a_{m,KM}+ u_{l,KM}b_{m,LM} \big) &= i\sum_k g_{klm}c_{k,KL}
.\label{JfQQB}
\end{align}
From these conditions, introducing the generator $Z^{11}_{KL}:=\sum_l c_{l,ML}B^{11}_l\in \g_{(1,1)}$, it becomes the $\Z_2^2$-central charge, which means
\begin{alignat}{2}
 [Q^{01}_{\ua\,K},\, Q^{10}_{\ub\,L}]&=\varepsilon_{\ua\ub} Z^{11}_{KL},\quad &
\llbracket Z^{11}_{KL}, \,T \rrbracket &=0
,\nn\\
 [Q^{01}_{\da\,K},\, Q^{10}_{\db\,L}]&=\varepsilon_{\da\db} \bar Z^{11}_{KL},&
\llbracket \bar Z^{11}_{KL}, \,T \rrbracket &=0,
\end{alignat} 
where $T\in \{ P_\mu, Z^{(1)}_{MN}, \bar Z^{(1)}_{MN}, Z^{(2)}_{MN}, \bar Z^{(2)}_{MN}, Z^{11}_{MN}, \bar Z^{11}_{MN}, Q_{\ua\,K}^{01}, \bar Q_{\da\,K}^{01}, Q_{\ua\,K}^{10}, \bar Q_{\da\,K}^{10}\}$.
From these observations, we can obtain the novel possible generators of symmetries of the $S$-matrix:
\begin{alignat}{2}
 \g_{(0,0)}&\ni P_\mu, M_{\mu\nu}, B^{00}_l, Z^{(1)}_{KL}
, \bar Z^{(1)}_{KL}, Z^{(2)}_{KL}, \bar Z^{(2)}_{KL}
,\qquad&
 \g_{(1,1)}&\ni B^{11}_l, Z^{11}_{KL}, \bar Z^{11}_{KL}
,\nn\\
 \g_{(0,1)}&\ni Q_{\ua\,K}^{01}, \bar Q_{\da\,K}^{01}
,\qquad&
 \g_{(1,0)}&\ni Q_{\ua\,K}^{10}, \bar Q_{\da\,K}^{10}
.
\end{alignat}
Therefore the following theorem holds.
\begin{thm}
There exits the unique $\Z_2^2$-graded extension of supersymmetric algebra which is given by the $\Z_2^2$-graded LSA $\langle P_\mu, Z^{(1)}_{KL}, \bar Z^{(1)}_{KL}, Z^{(2)}_{KL}, \bar Z^{(2)}_{KL}, Z^{11}_{KL}, \bar Z^{11}_{KL}, Q_{\ua\,K}^{01}, \bar Q_{\da\,K}^{01}, Q_{\ua\,K}^{10}, \bar Q_{\da\,K}^{10}\rangle$. We call this $\mathcal{N}=(K|L)$ $\Z_2^2$-supersymmetric algebra, where $K$ and $L$ are the number of the $(0,1)$ and $(1,0)$-graded supercharges.
\end{thm}
Specifically, if we set the number of each supercharges $Q_{\ua\,K}^{01}, Q_{\ua\,L}^{10}$ to $K=L=1$, it corresponds to the $\mathcal{N}=2\ \Z_2^2$-supersymmetric algebra, which was introduced by Bruce\cite{Bruce}.  
In the present case of $\Z_2^2$-grading, the internal symmetries are generated by ordinary Lie algebra $B^{00}_l$ and $B^{11}_l$. This does not seem to be a big difference from the ordinary supersymmetry. However we will see that internal symmetries are generated by a graded Lie algebra when $\Z_2^n$-graded ($n\geq 3$) extensions of supersymmetry are considered.
\begin{comment}
 By changing our convention to {\color{red} HOW?}, our algebra is completely consistent with the algebra in \cite{Bruce} which has Majorana fermionic supercharges with the $\Z_2^2$-gradings.
\end{comment}
%%%%%%%%%%%%%%%%%%%%%%%%%%%%%%%%%%%%%%%%%%%%%%%%%%%%%%%%%%%%%%%%%%%%%%%%%
\subsection{Possible $\Z_2^n$-graded generators of symmetries}
%%%%%%%%%%%%%%%%%%%%%%%%%%%%%%%%%%%%%%%%%%%%%%%%%%%%%%%%%%%%%%%%%%%%%%%%%
The discussion for the $\Z_2^2$-graded LSA in the previous section can also be applied to the $\Z_2^n$-graded LSA. We define two subsets of $\Z_2^n$:
\begin{align}
 (\Z_2^n)_{\ev}&:=\{\vec{b}\in \Z_2^n\  |\ \vec{b}\cdot \vec{b}=0 \mod 2\}\nn\\
 (\Z_2^n)_{\od}&:=\{\vec{f}\in \Z_2^n\  |\ \vec{f}\cdot \vec{f}=1 \mod 2\}
,
\end{align}
and the subalgebras of \eqref{Z2n-VS}:
\begin{align}
 \h_{\vec{b}}=\g_{\vec{0}}\oplus \g_{\vec{b}}\quad (\vec{b}\ne \vec{0})
,\qquad
 \h_{\vec{f}}=\g_{\vec{0}}\oplus \g_{\vec{f}}
,
\end{align}
where $\vec{0}$ is the identity element of $\Z_2^n$. 
Then $\h_{\vec{b}}$ is a Lie algebra, while $\h_{\vec{f}}$ is a LSA. This shows that we can obtain generators by applying the \CM theorem to $\h_{\vec{b}}$ and the \HLS theorem to $\h_{\vec{f}}$.  It can be seen that the following generators of the $S$-matrix are allowed:
\begin{alignat}{3}
 \g_{\vec{0}}&\ni P_\mu, M_{\mu\nu}, B^{\vec{0}}_l
,\qquad&
 \g_{\vec{b}}&\ni B^{\vec{b}}_l
,\qquad&
 \g_{\vec{f}}&\ni Q_{\ua\,K}^{\vec{f}}, \bar Q_{\da\,K}^{\vec{f}}
,
\end{alignat}
and the algebraic relations become
\begin{alignat}{2}
\llbracket Q_{\ua\,K}^{\vec{f}}, \bar Q_{\da\,L}^{\vec{f}\,'} \rrbracket &=\sigma_{\ua\da}^{\mu} P_\mu\delta_{KL}\delta^{\vec{f}\vec{f}\,'}
, \qquad&
\llbracket Q_{\ua\,K}^{\vec{f}}, Q_{\ub\,L}^{\vec{f}\,'} \rrbracket
&=\varepsilon_{\ua\ub} \sum_l a^{(\vec{f},\vec{f}\,')}_{l,KL} B^{\vec{f}+\vec{f}\,'}_l
,\nn\\
\llbracket Q^{\vec{f}}_{\ua\,L},B_l^{\vec{b}} \rrbracket 
&=s^{(\vec{f},\vec{b})}_{l,LM}Q^{\vec{b}+\vec{f}}_{\ua\,M}
, &
\llbracket B_l^{\vec{b}},B_m^{\vec{b}\,'} \rrbracket 
&=i\sum_k h^{(\vec{b},\vec{b}\,')}_{lmk} B_k^{\vec{b}+\vec{b}\,'}
,\nn\\
\llbracket \bar Q_{\da\,K}^{\vec{f}}, \bar Q_{\db\,L}^{\vec{f}\,'} \rrbracket
&=-(-1)^{(\vec{f},\vec{f}\,')}\varepsilon_{\da\db} \sum_l &\bar a^{(\vec{f},\vec{f}\,')}_{l,KL} B^{\vec{f}+\vec{f}\,'}_l
,
\label{Z2n-SUSY}
\end{alignat}
where $a,\,s$, and $h$ are structure constants with the $\vec{0}$-grading, and $a^{(\vec{f},\vec{f}\,')}_{l,KL}=(-1)^{\vec{f} \cdot \vec{f}\,' }a^{(\vec{f}\,',\vec{f})}_{l,LK}$ and $h^{(\vec{b},\vec{b}\,')}_{lmk}=-(-1)^{\vec{b}\cdot\vec{b}\,'}h^{(\vec{b}\,',\vec{b})}_{mlk}$.

In the same way as $\Z_2^2$-extension, we obtain  the conditions for structure constants from the $\Z_2^n$-Jacobi identities \eqref{Z2n-JI}. They are obtained from $(A,B,C)=(Q,\bar Q, B)$:
\begin{align}
 s_{l,LM}^{(\vec{f},\vec{f}+\vec{f}\,')}
&= \bar s_{l,ML}^{(\vec{f}\,',\vec{f}+\vec{f}\,')}
,
\end{align}
from $(A,B,C)=(Q, Q,\bar Q)$:
\begin{align}
 \sum_{l} a^{(\vec{f},\vec{f}\,')}_{l,LM} \bar s^{(\vec{f}\,'',\vec{f}+\vec{f}\,')}_{l,NK} &= 0 
,
\end{align}
from $(A,B,C)=(Q, Q, Q)$:
\begin{align}
 \sum_{l} a^{(\vec{f}\,',\vec{f}\,'')}_{l,LM} s^{(\vec{f},\vec{f}\,'+\vec{f}\,'')}_{l,KN} -(-1)^{\vec{f}\cdot\vec{f}\,'} \sum_l a^{(\vec{f},\vec{f}\,'')}_{l,KM} s^{(\vec{f}\,',\vec{f}+\vec{f}\,'')}_{l,LN} 
&= 0 
,
\end{align}
and from $(A,B,C)=(Q, Q, B)$:
\begin{align}
 s^{(\vec{f}\,',\vec{b})}_{l,LM}a^{(\vec{f},\vec{f}\,'+\vec{b})}_{m,KM}+(-1)^{\vec{f}\cdot\vec{f}\,'}s^{(\vec{f},\vec{b})}_{l,KM}a^{(\vec{f}\,',\vec{f}+\vec{b})}_{m,LM}&= \sum_{k} ih^{(\vec{f}+\vec{f}\,',\vec{b})}_{klm}a^{(\vec{f},\vec{f}\,')}_{k,KL}
.
\end{align}
Using these conditions and setting $Z_{KL}^{\vec{f}+\vec{f}\,'}:=\sum_l a^{(\vec{f},\vec{f}\,')}_{l,KL} B^{\vec{f}+\vec{f}\,'}_l$, then it satisfies
\begin{align}
 \llbracket Z_{KL}^{\vec{f}+\vec{f}\,'} ,\, T \rrbracket =
 \llbracket \bar Z_{KL}^{\vec{f}+\vec{f}\,'} ,\, T \rrbracket =0
,\quad
T\in\{P_\mu, Z_{KL}^{\vec{f}+\vec{f}\,'}, \bar Z_{KL}^{\vec{f}+\vec{f}\,'}, Q_{\ua\, K}^{\vec{f}}, \bar Q_{\ua\, K}^{\vec{f}}\},
\end{align}
where $Z_{KL}^{\vec{f}+\vec{f}\,'}=(-1)^{\vec{f}\cdot\vec{f}\,'}Z_{LK}^{\vec{f}+\vec{f}\,'}$.

Therefore the following theorem holds.
\begin{thm}
There exits the unique $\Z_2^n$-graded extension of supersymmetric algebra which is given by the $\Z_2^n$-graded LSA $\langle P_\mu, Z_{KL}^{\vec{f}+\vec{f}\,'}, \bar Z_{KL}^{\vec{f}+\vec{f}\,'}, Q_{\ua\, K}^{\vec{f}}, \bar Q_{\ua\, K}^{\vec{f}} \rangle$. 
\end{thm}
This algebra contains $K_i$ supercharges $Q_{\ua\, K_i}^{\vec{f}}$ in each subspace $\g_{\vec{f}}$ and nontrivial centers in each subspace $\mathfrak{g}_{\vec{b}}$. We call this $\mathcal{N}=(K_1| K_2| ... |K_{2^{n-1}})$ $\Z_2^n$-supersymmetric algebra.
Similarly to the $\Z_2^2$-case, if we set the number of each $Q_{\ua, K}^{\vec{f}}$ to $K_i=1$, it corresponds to the $\mathcal{N}=2^{n-1}\ \Z_2^n$-supersymmetric algebra, which was introduced by Bruce\cite{Bruce}. 
%%%%%%%%%%%%%%%%%%%%%%%%%%%%%%%%%%%%%%%%%%%%%%%%%%%%%%%%%%%%%%%%%%%%%%%%%
\subsection{$\Z_2^n$-Graded internal symmetries}
In this section, we focus on the generators of internal symmetry with $\Z_2^n$-grading $B_l^{\vec{b}}$. The discussion in the previous section says there can exist $B_l^{\vec{b}}$ in each subspace $\g_{\vec{b}}$ and they satisfy 
\begin{align}
 \llbracket B_l^{\vec{b}},B_m^{\vec{b}\,'} \rrbracket 
&=i\sum_k h^{(\vec{b},\vec{b}\,')}_{lmk} B_k^{\vec{b}+\vec{b}\,'}.
\tag{\ref{Z2n-SUSY}}
\end{align}
This means the set of $B_{l}^{\vec{b}}$ is a $(\Z_2^n)_{\ev}$-graded LSA. Especially, if we set $n=2k+1$, we can prove that this algebra is isomorphic to the $\Z_2^{2k}$-graded Lie algebra, which is another type of $\Z_2^n$-graded extension of Lie algebra introduced by Rittenberg and Wyler\cite{RW1,RW2}. 

\begin{dfn}\label{defLA}
Let $\g_{LA}$ be a $\Z_2^{2n}$-graded vector space \eqref{Z2n-VS} and $\vec{\ua}=(\vec{\ua}_1,\vec{\ua_2},...,\vec{\ua}_n)\in \Z_2^{2n}$ where $\vec{\ua}_i=(\ua_{2i-1},\ua_{2i})\in \Z_2^2$. Then $\g_{LA}$ is called $\Z_2^{2n}$-graded Lie algebra if $\g_{LA}$ has the $\Z_2^{2n}$-Lie bracket defined by
\begin{align}
 [A,B\}:=AB-(-1)^{\sum_{i=1}^n |\vec{\ua}_i\times \vec{\ub}_i|}BA,\quad A\in \g_{\vec{\ua}},\  B\in \g_{\vec{\ub}}
,\label{Z2n-LieB}
\end{align}
and $[A,B\}$ satisfies 
\begin{align}
 [A,B\}&\in\g_{\vec{\ua}+\vec{\ub}},\nn\\
 [A,[B,C\}\}&=[[A,B\},C\}+(-1)^{\sum_{i=1}^n |\vec{\ua}_i\times \vec{\ub}_i|}[B[A,C\}\}
\end{align}
\end{dfn}

\begin{thm}
 The $(\Z_2^{2n+1})_{\ev}$-graded LSA 
\begin{align}
\g_{\ev}:=\bigoplus_{\vec{b}\in(\Z_2^{2n+1})_{\ev}}\g_{\vec{b}}
\end{align}
 is isomorphic to the $\Z_2^{2n}$-graded Lie algebra $\g_{LA}$.
\end{thm}
\begin{proof}
Let $f$ be an injective map $f: \vec{\ua}\in \Z_2^{2n}\to\vec{a}\in(\Z_2^{2n+1})_{\ev}$, defined by
\begin{align*}
\vec{a}=\left( \vec{\ua}_1, \vec{\ua}_2+|\vec{\ua}_1|(1,1),\cdots \vec{\ua}_i+\sum_{k=1}^{i-1}|\vec{\ua}_k|(1,1),\cdots, \vec{\ua}_{n}+\sum_{k=1}^{n-1}|\vec{\ua}_k|(1,1), \sum_{k=1}^{n}|\vec{\alpha}_k| \right)
\end{align*}
where  $|\vec{a}_k|^2:=a_{2k-1}^2+a_{2k}^2$, and it is equal to $|\vec{a}_k|=a_{2k-1}+a_{2k} \mod 2$.

Since $\Z_2^{2n}$-graded Lie algebra $\g_{LA}$ and $(\Z_2^{2n+1})_{\ev}$-graded LSA $\g_{\ev}$ have the same number of the graded subspaces, the only difference of them is the Lie bracket, which comes from $\sum_{i=1}^n |\vec{\ua}_i\times \vec{\ub}_i|$ in \eqref{Z2n-LieB} and $\vec{a}\cdot \vec{b}$ in \eqref{Z2n-LB}. However, taking into account that computations are performed in modulo 2, we have straightforwardly 
\begin{align}
 \sum_{i=1}^n |\vec{\ua}_i\times \vec{\ub}_i|=\vec{a}\cdot \vec{b}.
\end{align}
Therefore the map $f$ is an isomorphism between $\g_{LA}$ and $\g_{\ev}$.
\end{proof}
From this theorem, we can redefine $\Z_2^{2n}$-graded Lie algebra (Definition \ref{defLA}).
\begin{dfn}
 Let $\g$ be a $\Z_2^{n}$-graded LSA in Definition \ref{defsalg}. We define $\Z_2^{n-1}$-graded Lie algebra as the subalgebra 
\begin{align}
 \g_{\ev}:=\bigoplus_{\vec{b}\in(\Z_2^n)_{\ev}}\g_{\vec{b}}\subset \g.
\end{align}
\end{dfn}
This is a more general definition, which means we can define not only $\Z_2^{2n}$-graded Lie algebras, but also $\Z_2^{2n+1}$-graded Lie algebras. Finally, for the internal symmetry, the following theorem holds.
\begin{thm}
 The $\Z_2^n$-graded Lie algebra is allowed as the generators of the internal symmetry of S-matrix.
\end{thm}
This is very interesting result because we can consider new types of gauge symmetries in relativistic QFT.
Specifically, $\Z_2^2$-extension of $\mathfrak{sl}_n$ has already been introduced in \cite{RW2} and some representation of $\Z_2^2$-$\mathfrak{sl}_2$ was investigated in \cite{AIKTT}. These will be helpful to construct $\Z_2^2$-$\mathfrak{sl}_n$ gauge field theory.
%%%%%%%%%%%%%%%%%%%%%%%%%%%%%%%%%%%%%%%%%%%%%%%%%%%%%%%%%%%%%%%%%%%%%%%%%

%%%%%%%%%%%%%%%%%%%%%%%%%%%%%%%%%%%%%%%%%%%%%%%%%%%%%%%%%%%%%%%%%%%%%%%%%
\section{Concluding remarks}
In this paper, we explored novel possible generators of symmetries of $S$-matrix in order to apply $\Z_2^n$-graded Lie superalgebra to relativistic quantum field theory. As a result, we found $\mathcal{N}=(K_1|K_2|\cdots |K_{2^{n-1}})$ $\Z_2^n$-graded supersymmetry is the only allowed extension of supersymmetric algebra as the $\Z_2^n$-graded extension. Specifically, if we set the number of each $Q_{\ua\, K}^{\vec{f}}$'s to $K_i=1$, they correspond to the $\mathcal{N}=2^{n-1}\ \Z_2^n$-supersymmetric algebra, which was introduced by Bruce\cite{Bruce}. However we cannot determine the number of supercharges purely from algebraic restrictions. This should be determined by the relationship between the degrees of freedom of particles within a model under consideration. 

Nontrivial results of the present paper are summarized as follows.
\begin{enumerate}
 \item It is possible to construct $\Z_2^n$-graded supersymmetric relativistic QFT.
 \item Gauge symmetry is not necessary generated by a Lie algebra, but by a $\Z_2^n$-graded Lie algebra.   
\end{enumerate} 
Thus an important future work is to construct a $\Z_2^2$-$\mathfrak{sl}_2$ gauge field theory.

Finally, in this paper, we have only considered the case of massive particles.
According to \cite{CM,HLS}, conformal generators can exist in the case of massless particles.
Therefore if we explore novel generators in the massless case, we could obtain more nontrivial $\Z_2^n$-graded conformal generators.  

\section*{Acknowledgment}
This work was supported by JST SPRING, Grant Number JPMJSP2139 (RI and AN). RI and AN would like to express our sincere gratitude to Professor Naruhiko Aizawa and Professor Nobuhito Maru, at Osaka Metropolitan University, for their invaluable guidance and support throughout our research.


%%%%%%%%%%%%%%%%%%%%%%%%%%%%%%%%%%%%%%%%%%%%%%%%%%%%%%%%%%%%%%%%%%%%%%%%%%
\begin{thebibliography}{99}
%%%%%%%%%%%%%%%%%%%%%%%%%%%%%%%
%\bibitem{E.W} E. Wigner, \textit{On the Consequences of the Symmetry of the Nuclear Hamiltonian on the Spectroscopy of Nuclei}, Phys. Rev. {\bf 51}, 106 (1937),

\bibitem{F.J.D} F. J. Dyson, \textit{Symmetry Groups in Nuclear and particle Physics}, W. A. Benjamin, New York (1966),

\bibitem{S.W} S. Weinberg, \textit{The Quantum Theory of Fields:Volume III: Supersymmetry}, Cambridge University Press (2000),

\bibitem{M.G} M. Gell-Mann, \textit{A schematic model of baryons and mesons}, Phys. Lett. {\bf 8}, 3 (1964),

\bibitem{S.C} S. Coleman, \textit{Trouble with Relativistic SU(6)}, Phys. Rev. {\bf 138}, B1262 (1965),

\bibitem{CM} S. Coleman and J. Mandula, \textit{All Possible Symmetry of S-matrix}, Phys. Rev. {\bf 159}, 1251 (1967),

\bibitem{J.M} J. Mandula, \textit{Coleman-Mandula theorem}, Scholapedia, 10(2):7476 (2015),

\bibitem{HLS} R. Haag, J. T. \L opusza\'nski, M. Sohnius,
\textit{All possible generators of supersymmetry of the S-matrix}, 
Nucl. Phys. B {\bf 88}, 257 (1975),
%%%%%%%%%%%%%%%%%%%%%%%%%%%%%%%%%%%%%%%%%%%%%%%%%%%%%%%%%%%%%%%%%%%%

\bibitem{RW1} V. Rittenberg and D. Wyler, {\it  Generalized superalgebras}, Nucl. Phys. B {\bf 139}, 189 (1978).

\bibitem{RW2} V. Rittenberg and D. Wyler, {\it Sequences of $\mathbb{Z}_2 \otimes \mathbb{Z}_2$ graded Lie algebras and superalgebras}, 
J. Math. Phys. {\bf 19}, 2193  (1978).

\bibitem{Ree} R. Ree, 
\textit{Generalized Lie elements}, Canad. J. Math. \textbf{12}, 493 (1960).

\bibitem{sch1} M. Scheunert, \textit{Generalized Lie algebras}, J. Math. Phys. {\bf 20}, 712 (1979).

%%%%%%%%%%%%%%%%%%%%%%%%%%%%%%%%%%%%%%%%%%%%%%%%%%%%%%%%%%%%%%
\bibitem{Vas} M . A. Vasiliev, \textit{de Sitter supergravity with positive cosmological constant and generalised Lie superalgebras}, Class. Quantum Gravity \textbf{2}, 645 (1985).

\bibitem{jyw} P. D. Jarvis, M. Yang and B. G. Wybourne, 
\textit{Generalized quasispin for supergroups},  J. Math. Phys. {\bf 28},  1192 (1987).

\bibitem{tol} V. N. Tolstoy, 
\textit{Once more on parastatistics}, Phys. Part. Nucl. Lett. {\bf{11}},  933 (2014).

\bibitem{AKTT1} N. Aizawa,  Z. Kuznetsova, H. Tanaka, F. Toppan,
\textit{${\mathbb Z}_2\times {\mathbb Z}_2$-graded Lie Symmetries of the L\'evy-Leblond Equations}, 
Prog. Theor. Exp. Phys. \textbf{2016}, 123A01 (2016).

\bibitem{AKTT2} N. Aizawa,  Z. Kuznetsova, H. Tanaka, F. Toppan,
\textit{Generalized supersymmetry and L\'evy-Leblond equation}, 
in ``Physical and Mathematical Aspects of Symmetries,''
J.-P. Gazeau, S. Faci, T. Micklitz, R. Scherer, F. Toppan (editors),  
Springer (2017) p.79.

%%%%%%%%%%%%%%%%%%%%%%%%%%%%%%%%%%%%%%%%%%%%%%%%%%%%%%%%%%%%%%%%%%%%%
\bibitem{zhe} A. A. Zheltukhin, 
\textit{Para-Grassmann extension of the Neveu-Schwartz-Ramond algebra}, Theor. Math. Phys. {\bf 71},  491  (1987) (Teor. Mat. Fiz. {\bf 71}  218 (1987)).

\bibitem{LR} B. Le Roy, 
\textit{$\mathbb{Z}_n^3$-Graded colored supersymmetry}, Czech. J. Phys. \textbf{47}, 47 (1997).

\bibitem{Toro1} L. A. Wills-Toro, 
\textit{$(I,q)$-graded Lie algebraic extensions of the Poincar\'e algebra, constraints on $I$ and $q$}, J. Math. Phys. \textbf{36}, 2085 (1995).

\bibitem{Toro2} L. A. Wills-Toro, 
\textit{Trefoil symmetries I. Clover extensions beyond Coleman-Mandula theorem}, J. Math. Phys. \textbf{42}, 3915  (2001). 

\bibitem{tol2} V. N. Tolstoy,  
\textit{Super-de Sitter and Alternative Super-Poincar\'e Symmetries},  
In: Dobrev V. (eds) Lie Theory and Its Applications in Physics. Springer Proceedings in Mathematics \& Statistics, 
vol. 111, Springer, Tokyo, 2014.

\bibitem{Bruce} A. J. Bruce, 
\textit{On a $\mathbb{Z}_2^n$-graded version of supersymmetry}, Symmetry \textbf{11}, 116 (2019).

%%%%%%%%%%%%%%%%%%%%%%%%%%%%%%%%%%%%%%%%%%%%%%%%%%%%%%%%%%%%%

\bibitem{BruDup} A. J. Bruce and S. Duplij, 
{\it Double-graded supersymmetric quantum mechanics},
J. Math. Phys.  \textbf{61}, 063503 (2020).

\bibitem{AAD}N. Aizawa, K. Amakawa, S. Doi, 
{\it $\mathcal{N}$-Extension of double-graded supersymmetric and superconformal quantum mechanics}, 
J. Phys. A: Math. Theor. \textbf{53}, 065205 (2020).


\bibitem{AAd2} N. Aizawa, K. Amakawa, S. Doi, 
{\it $\mathbb{Z}_2^n$-Graded extensions of supersymmetric quantum mechanics via Clifford algebras}, 
J. Math. Phys. \textbf{61}, 052105 (2020).


\bibitem{AKTcl} N. Aizawa, Z. Kuznetsova and F. Toppan,
{\it ${\mathbb Z}_2\times {\mathbb Z}_2$-graded mechanics: the classical theory},  
Eur. Phys. J. C \textbf{80}, 668 (2020).

\bibitem{AKTqu} N. Aizawa, Z. Kuznetsova and F. Toppan,
{\it ${\mathbb Z}_2\times {\mathbb Z}_2$-graded mechanics: the quantization},  
Nucl. Phys. B \textbf{967},  115426 (2021).


\bibitem{DoiAi1} S. Doi and N. Aizawa,
\textit{$\mathbb{Z}_2^3$-Graded extensions of Lie superalgebras and superconformal quantum mechanics}, 
SIGMA	\textbf{17} 071, (2021).

\bibitem{DoiAi2} S. Doi and N. Aizawa,
\textit{Comments of $\mathbb{Z}_2^2$-supersymmetry in superfield formalism}, 
Nucl. Phys.  \textbf{B974}, 115641 (2022).

\bibitem{AiDoi} N. Aizawa and S. Doi, 
\textit{Irreducible representations of $\mathbb{Z}_2^2$-graded $\mathcal{N} = 2$ supersymmetry algebra and $\mathbb{Z}_2^2$-graded supermechanics}, 
J. Math. Phys. \textbf{63}, 091704 (2022).

\bibitem{brusigma} A. J. Bruce, \textit{$\mathbb{Z}_2 \times \mathbb{Z}_2$-graded supersymmetry: 2-d sigma models},
J. Phys. A:Math. Theor. \textbf{53}, 455201 (2020).

\bibitem{bruSG} A. J. Bruce, \textit{Is the $\mathbb{Z}_2 \times \mathbb{Z}_2$-graded sine-Gordon equation integrable ?}, 
Nucl. Phys. B \textbf{971}, 115514 (2021).

\bibitem{AIKT} N. Aizawa, R. Ito, Z. Kuznetsova, F. Toppan, 
\textit{New aspects of the $\mathbb{Z}_2\times\mathbb{Z}_2$-graded 1D superspace: induced strings and 2D relativistic models}, 
Nucl. Phys. B \textbf{991},  116202 (2023).


\bibitem{Topp} F. Toppan, {\it $\mathbb{Z}_{2} \times \mathbb{Z}_{2}$-graded parastatics in multiparticle quantum Hamiltonians}, J. Phys. A: Math. Theor. \textbf{54}, 115203 (2021).
%arXiv:2008.11554 [hep-th].

\bibitem{Topp2} F. Toppan, \textit{Inequivalent quantizations from gradings and $ \mathbb{Z}_2 \times \mathbb{Z}_2$ parabosons}, 
J. Phys. A, Math. Theor. \textbf{54}, 355202 (2021).
%	arXiv:2104.09692 [hep-th].

\bibitem{AiItoTa}
N. Aizawa, R. Ito and T. Tanaka, {\textit{$\mathbb{Z}_2^2$-graded supersymmetry via superfield on minimal 
$\mathbb{Z}_2^2$-superspace}}, 
J. Phys. A: Math. Theor. {\bf 57}  435201 (2024).
\bibitem{AiItoTa2}
N. Aizawa, R. Ito and T. Tanaka, {\textit{$\mathcal{N}=2$ Double-graded supersymmetric quantum mechanics via dimensional reduction}},
AIMS Math. {\bf 9} 10494 (2024).

%%%%%%%%%%%%%%%%%%%%%%%%%%%%%%%%%%%%%%%%%%%%%%%%%%%%%%%%%%%%%%%%%%%%%%%

\bibitem{sch3} M. Scheunert, \textit{Casimir elements of $\epsilon$-Lie algebras}, 
J. Math. Phys. \textbf{24}, 2671 (1983).

\bibitem{SchZha} M. Scheunert and R. B. Zhang, 
\textit{Cohomology of Lie superalgebras and their generalizations}, 
J. Math. Phys. \textbf{39}, 5024 (1998). 

\bibitem{Sil} S. D. Silvestrov, 
\textit{On the classification of 3-dimensional coloured Lie algebras},  
Banach Center Publications, \textbf{40}, 159 (1997).


 \bibitem{ChSiVO} X.-W. Chen, S. D. Silvestrov and F. Van Oystaeyen, 
\textit{Representations and cocycle twists of color Lie algebras}, 
Algebr. Represent. Theor. \textbf{9}, 633 (2006).

\bibitem{SigSil} G. Sigurdsson and S. D. Silvestrov, 
\textit{Bosonic realizations of the colour Heisenberg Lie algebras},
Nonlinear Math. Phys. \textbf{13}, supplement  110 (2006).

\bibitem{PionSil} D. Piontkovski and S. D. Silvestrov, 
\textit{Cohomology of 3-dimensional color Lie algebras},
J. Algebra \textbf{316}, 499 (2007).


\bibitem{CART} R. Campoamor-Stursberg and M. Rausch de Traubenberg, 
\textit{Color Lie algebras and Lie algebras of order $F$},
J. Gen. Lie Theory Appl. \textbf{3}, 113 (2009). 

\bibitem{MohSal} M. Mohammadi and H.  Salmasian, \textit{The Gelfand-Naimark-Segal construction for unitary representations of $\mathbb Z_2^n$-graded Lie supergroups}, 
Banach Cent. Publ. \textbf{113}, 263 (2017). 


\bibitem{NAJS} N. Aizawa and J. Segar, 
\textit{$\mathbb{Z}_2 \times \mathbb{Z}_2$ generalizations of ${\cal N} = 2$ super Schr\"odinger  algebras and their representations}, 
J. Math. Phys. \textbf{58}, 113501 (2017).

\bibitem{NAPIJS} N. Aizawa, P. S. Isaac and J. Segar,
\textit{${\mathbb Z}_2\times {\mathbb Z}_2$ generalizations of ${\cal N} = 1$  superconformal Galilei algebras and their representations}, 
J.  Math. Phys. \textbf{60}, 023507 (2019). 

\bibitem{StoVDJ} N. I. Stoilova, J. Van der Jeugt, 
\textit{The $\mathbb{Z}_2 \times \mathbb{Z}_2$-graded Lie superalgebra $ \mathfrak{pso}(2m+1|2n) $ and new parastatistics representations}, 
J. Phys. A:Math. Theor. \textbf{51}, 135201 (2018).

\bibitem{Meyer} P. Meyer, 
\textit{The Kostant invariant and special $\epsilon$-orthogonal representations for $\epsilon$--quadratic colour Lie algebras}, 
J. Algebra \textbf{572}, 337 (2021).


\bibitem{IsStvdJ} P. S. Isaac, N. I. Stoilova, J. van der Jeugt, 
\textit{The $\mathbb{Z}_2 \times \mathbb{Z}_2$-graded general linear Lie superalgebra},
J. Math. Phys. \textbf{61}, 011702 (2020).

\bibitem{NAPSIJSinf}
N. Aizawa, P. S. Isaac and J. Segar, 
\textit{$\mathbb{Z}_2 \times \mathbb{Z}_2$ generalizations of infinite dimensional Lie superalgebra of conformal type with complete classification of central extensions}, 
Rep. Maht. Phys.  \textbf{85}, 351 (2020).

\bibitem{NAKA} K. Amakawa and N. Aizawa, 
\textit{A classification of lowest weight irreducible modules over $\mathbb{Z}_2^2$-graded extension of $osp(1|2)$}, 
J. Math. Phys. \textbf{62}, 043502 (2021).

\bibitem{KTclassification} Z. Kuznetsova, F. Toppan, 
\textit{Classification of minimal $\mathbb{Z}_2 \times \mathbb{Z}_2$-graded Lie (super)algebras and some applications}, J. Math. Phys. \textbf{62}, 063512 (2021).


\bibitem{Meyer2} P. Meyer, 
\textit{Cubic Dirac operators and the strange Freudenthal-de Vries formula for colour Lie algebras},
Transform. Groups  \textbf{27}, 1307 (2022).

\bibitem{NeliJoris} N. I. Stoilova, J. van der Jeugt, 
\textit{The ${\mathbb{Z}}_{2}\times {\mathbb{Z}}_{2}$-graded Lie superalgebras $\mathfrak{p}\mathfrak{s}\mathfrak{o}(2n+1\vert 2n)$ and $\mathfrak{p}\mathfrak{s}\mathfrak{o}(\infty \vert \infty )$, and parastatistics Fock spaces}, 
J. Phys. \textbf{A:}Math. Theor. \textbf{55}, 045201 (2022).



\bibitem{NeliJoris2} N. I. Stoilova, J. van der Jeugt, 
\textit{On classical $\mathbb{Z}_2 \times \mathbb{Z}_2$-graded Lie algebras}, 
J. Math. Phys. \textbf{64}, 061702 (2023). 

\bibitem{LuTan} R. Lu and Y. Tan, 
\textit{Construction of color Lie algebras from homomorphisms of modules of Lie algebras}, 
J. Algebra \textbf{620}, 1 (2023).


\bibitem{FaFaJ} F. S. Alshammari, Md F. Hoque, J. Segar, 
\textit{The $osp(1|2)$ $ \mathbb{Z}_2 \times \mathbb{Z}_2$ graded algebra and its irreducible representations}, 
arXiv:2306.12381 [math-ph]. 

\bibitem{AiSe} N. Aizawa, J. Segar,
\textit{Affine extensions of $\mathbb{Z}_2^2$-graded $osp(1|2)$ and Virasoro algebra},
arXiv:2409.07938 [math-ph]

\bibitem{AIKTT}
N. Aizawa, R.Ito, Z. Kuznetsova, T. Tanaka, F, Toppan,
\textit{Integrable $\mathbb{Z}_2^2$-graded Extensions of the Liouville and Sinh-Gordon Theories},
arXiv:2406.13503 [math-ph].


\end{thebibliography}
\end{document}
